\PassOptionsToPackage{usenames,dvipsnames}{xcolor}
\documentclass[a4paper,12pt]{letter}
\usepackage[utf8]{inputenc}
\usepackage[T1]{fontenc}
\usepackage[english,slovene]{babel}
\usepackage[tikz]{bclogo}
\usepackage{qrcode,marvosym}
\usepackage{pgf,tikz,pgfplots}
\pgfplotsset{compat=1.14}
\usepackage{mathrsfs}
\usetikzlibrary{arrows}
\usepackage{graphicx}
\newcommand{\ds}{\displaystyle}
\usepackage{fancybox}
\usepackage{amsfonts}
\usepackage{amsmath}
\usepackage{wallpaper}
\usepackage{hyperref}
\usepackage{array}
\newcolumntype{M}[1]{>{\centering\arraybackslash}m{#1}}
\newcolumntype{N}{@{}m{0pt}@{}}
\usepackage{fancyhdr,enumerate}
\newcommand{\ora}{\overrightarrow}
\renewcommand{\headrulewidth}{3pt}
\renewcommand{\headrule}{\hbox to\headwidth{%
  \color{gray}\leaders\hrule height \headrulewidth\hfill}}
\renewcommand{\footrulewidth}{2pt}
\renewcommand{\footrule}{\hbox to\headwidth{%
  \color{brown}\leaders\hrule height \footrulewidth\hfill}}
\usepackage{gensymb}
\usepackage{amssymb}
  \textwidth 20 true cm
  \oddsidemargin -2 true cm
  \evensidemargin -2 true cm
  \topmargin -3.5 true cm
  \textheight 27.5 true cm
  \linespread{1.2}
\renewcommand{\headrulewidth}{0.3pt}
\renewcommand{\footrulewidth}{1.9pt}
\definecolor{qqwuqq}{rgb}{1,0,0}
\definecolor{qqqqff}{rgb}{0,0,1}
\definecolor{qqffqq}{rgb}{0,1,0}
\definecolor{ffqqqq}{rgb}{1,0,0}
\definecolor{zzuxux}{rgb}{0,0,0}
\usepackage{mdframed}
\usepackage{multicol}
\definecolor{bgblue}{RGB}{0,0,253}
\usepackage{eurosym}
%%%%%%%%%%%%%%%%%%%%%%%%%%%%%%%%%%%%%%%%%%%%%%%%%%%%%%%%
\newcounter{tocke}
\setcounter{tocke}{0}
\usepackage{tikz-3dplot}
%%%%%%%%%%%%%%%%%%%%%%%%%%%%%%%%%%%%%%%%%%%%%%%%%%%%%%%%%%
\mdfdefinestyle{theoremstyle}{%
linecolor=brown!40,linewidth=1pt,%
frametitlerule=true,%
frametitlebackgroundcolor=brown!50,
innertopmargin=\topskip,
}
\mdfdefinestyle{theoremstyle1}{%
linecolor=brown,middlelinewidth=1pt,%
frametitlerule=true,%
apptotikzsetting={\tikzset{mdfframetitlebackground/.append style={%
shade,left color=black!40, right color=gray!10},
mdfbackground/.append style={
top color=white!100!white,
bottom color=black!5
},
}},
frametitlerulecolor=gray,
frametitlerulewidth=2pt,
innertopmargin=\topskip,
innerbottommargin=\topskip,
}
\mdtheorem[style=theoremstyle1]{naloga}{Naloga}
\mdtheorem[style=theoremstyle]{resitev}{Analiza Naloge}
\mdtheorem[style=theoremstyle]{kriterij}{\v{S}tevilo dose\v{z}enih to\v{c}k na testu}
%%%%%%%%%%%%%%%%%%%%%%%%%%%%%%%%%%%%%%%%%%%%%%%%%%%%%%%%%%
\usepackage[table]{xcolor}
\usepackage[printwatermark]{xwatermark}
\newwatermark[allpages,color=brown!2,angle=45,scale=5,
xpos=-5,ypos=0]{matej.info}
%%%%%%%%%%%%%%%%%%%%%%%%%%%%%%%%%%%%%%%%%%%%%%%%%%%%%%%%%%

\begin{document}
\pagestyle{fancy}
\renewcommand\headrulewidth{0pt}
\lhead{}\chead{}\rhead{}
\rfoot{\vspace*{-2\baselineskip}}
\renewcommand{\vec}{\overrightarrow}
%%%%%%%%%%%%%%%%%%%% podatki testa -- tu spreminjaj
\def\programletnik{Gim - 4}
\def\naslov{Test 1.3.P - }
\def\sklop{Geometrijska telesa}
%%%%%%%%%%%%%%%%%%%% tocke po nalogah
\def\tockeA{4+3}
\def\tockeB{3+3+4}
\def\tockeC{3+4+3}
\def\tockeD{4+4}
%%%%%%%%%%%%%%%%%%%%
\begin{minipage}{20cm}
\begin{bclogo}[couleur=brown!30,
  arrondi=0,ombre=true, couleurOmbre=gray!50,
logo =\bcplume,
  noborder=true]
{\textrm{{\Large{\naslov \sklop} \Huge{}
\large{} \hfill $\programletnik$}}}
\end{bclogo}
\end{minipage}
\begin{minipage}{20cm}
\begin{bclogo}[couleur=brown!30,
  arrondi=0,ombre=true, couleurOmbre=gray!50,
logo =\bcplume,
  noborder=true]
{\textrm{{\raisebox{-0.3cm}{\large{Ime in Priimek: \rule{10cm}{0.1mm}} \Huge{}
\large{} \hfill $ $}}}}
\end{bclogo}
\end{minipage}
%%%%%%%%%%%%%%%%%%%%%%%%%%%%%%%%%%%%%%%%%%%%%%%%%

%%%%%%%%%%%%%%%%%%%%%%%%%%%%%%%%%%%%%%%%%%%%%%%%%
\begin{naloga}[\hfill $\tockeA$] \addtocounter{tocke}{\numexpr\tockeA\relax}
Pravilni \v{s}esterokotni prizmi z osnovnim robom $a=4$ cm in vi\v{s}ino $v=10$ cm je vrisan valj (dotika se vseh \v{s}estih stranskih ploskev).
\begin{enumerate}[a)]
\item Izra\v{c}unaj polmer vrisanega valja in njegovo prostornino.
\item Koliko odstotkov prostornine prizme zavzema valj?
\end{enumerate}
\end{naloga}\newpage

%%%%%%%%%%%%%%%%%%%%%%%%%%%%%%%%%%%%%%%%%%%%%%%%%
\begin{naloga}[\hfill $\tockeB$] \addtocounter{tocke}{\numexpr\tockeB\relax}
Pravilni tristrani piramidi je rob osnovne ploskve $a=6$ cm, stranski rob pa $b=4$ cm.
\begin{enumerate}[a)]
\item Izra\v{c}unaj vi\v{s}ino piramide.
\item Izra\v{c}unaj prostornino piramide.
\item Izra\v{c}unaj kot med stranskim robom in osnovno ploskvijo.
\end{enumerate}
\end{naloga}\newpage

%%%%%%%%%%%%%%%%%%%%%%%%%%%%%%%%%%%%%%%%%%%%%%%%%
\begin{naloga}[\hfill $\tockeC$] \addtocounter{tocke}{\numexpr\tockeC\relax}
Osni presek sto\v{z}ca je enakostrani\v{c}ni trikotnik s stranico $10$ cm.
\begin{enumerate}[a)]
\item Izra\v{c}unaj polmer in vi\v{s}ino sto\v{z}ca.
\item Izra\v{c}unaj prostornino in povr\v{s}ino sto\v{z}ca.
\item Izra\v{c}unaj kot med stranico sto\v{z}ca in osnovno ploskvijo.
\end{enumerate}
\end{naloga}\newpage

%%%%%%%%%%%%%%%%%%%%%%%%%%%%%%%%%%%%%%%%%%%%%%%%%
\begin{naloga}[\hfill $\tockeD$] \addtocounter{tocke}{\numexpr\tockeD\relax}
Lik, sestavljen iz pravokotnika s stranicama $8$ cm in $4$ cm ter polkroga s polmerom $4$ cm, ki je pripet na kraj\v{s}o stranico pravokotnika, zavrtimo okoli dalj\v{s}e stranice (osi simetrije) za polni kot.
\begin{enumerate}[a)]
\item Izra\v{c}unaj prostornino nastalega telesa.
\item Izra\v{c}unaj povr\v{s}ino nastalega telesa.
\end{enumerate}
\end{naloga}

\newpage
\vfill
\begin{kriterij*} \begin{center} \renewcommand{\arraystretch}{2.5}
\begin{tabular}{|>{\columncolor{gray!20}}c||c|c|c|c|c||c|c|c|} \hline Odstotki &$[0,45)$ & $[45,60)$ & $[60,75)$ & $[75,90)$ & $[90,100]$& \v{S}t. mo\v{z}nih to\v{c}k& Dose\v{z}ene to\v{c}ke & OCENA \\ \hline Ocena & 1 & 2 & 3 & 4 & 5 &\cellcolor{gray!20} $\thetocke$&\cellcolor{gray!20} \phantom{ocena} & \cellcolor{gray!20} \phantom{a} \\ \hline \end{tabular} \end{center} \vspace{0.3cm}  \end{kriterij*}

\newpage
\textsc{Re\v{s}itve.}

\begin{enumerate}
\item a) $r=2\sqrt3\ \text{cm}\approx3{,}46$ cm, $V_{\text{valj}}=120\pi\ \text{cm}^3\approx376{,}99\ \text{cm}^3$.\quad b) $V_{\text{prizma}}=240\sqrt3\ \text{cm}^3\approx415{,}69\ \text{cm}^3$; valj zavzema $\approx90{,}69\,\%$ prostornine prizme.
\item a) $v=2$ cm.\quad b) $V=6\sqrt3\ \text{cm}^3\approx10{,}39\ \text{cm}^3$.\quad c) $\varphi=30^{\circ}$.
\item a) $r=5$ cm, $v=5\sqrt3\ \text{cm}\approx8{,}66$ cm.\quad b) $V=\dfrac{125\sqrt3}{3}\pi\ \text{cm}^3\approx226{,}72\ \text{cm}^3$, $S=75\pi\ \text{cm}^2\approx235{,}62\ \text{cm}^2$.\quad c) $\varphi=60^{\circ}$.
\item a) $V=\dfrac{512\pi}{3}\ \text{cm}^3\approx536{,}17\ \text{cm}^3$ ($V_{\text{valj}}=128\pi$, $V_{\text{polkrogla}}=\tfrac{128\pi}{3}$).\quad b) $S=112\pi\ \text{cm}^2\approx351{,}86\ \text{cm}^2$.
\end{enumerate}
\end{document}
